%% kapa: exam=1 round=19 year=2008 period=1 form=A subject=appraisal_law subject_ko=부동산관계법규 seq=5 number=45
%% source: Q-Net 감정평가사 기출문제 (cst00309) artlSeq=1001962 file=extracted/1008/1001962/2008년 제19회 감정평가사 1차시험 기출문제/1교시.hwp
%% conversion: kordoc 4.24.51 -> tools/gen1.py (auto); review=proofread-sonnet; flags=none
\begin{question}{45}
국토의 계획 및 이용에 관한 법령상 일정기간 개발행위허가를 제한 또는 그 제한의 연장을 할 수 있는 지역이 아닌 것은?
\begin{choices}
  \item 녹지지역으로 수목이 집단적으로 생육되고 있거나 조수류 등이 집단적으로 서식하고 있는 지역
  \item 개발행위로 인하여 주변환경 등이 크게 오염되거나 손상의 우려가 있는 지역
  \item 지구단위계획구역으로 지정된 지역
  \item 도시관리계획을 수립하고 있는 지역으로서 당해 도시관리계획이 결정될 경우 용도지역의 변경이 예상되고 그에 따라 개발행위허가의 기준이 크게 달라질 것으로 예상되는 지역
  \item 개발제한구역으로 지정된 지역
\end{choices}
\end{question}
\answer{5}
